\documentclass[times,twocolumn,final]{elsarticle}

\usepackage{cag}
\usepackage{framed,multirow}
\usepackage{amssymb}
\usepackage{amsmath}
\usepackage[switch]{lineno}
\usepackage{graphicx}
\usepackage{booktabs}
\usepackage{url}
\usepackage{placeins}
\setlength{\emergencystretch}{2em}
\setcounter{totalnumber}{6}
\setcounter{topnumber}{4}
\setcounter{bottomnumber}{4}
\renewcommand{\topfraction}{0.92}
\renewcommand{\dbltopfraction}{0.95}
\renewcommand{\textfraction}{0.05}
\renewcommand{\floatpagefraction}{0.85}


\journal{Computers \& Graphics}

\begin{document}

\begin{frontmatter}

\title{Control Allocation in Multi-view Diffusion Texturing: An Empirical Evaluation of TCAS and Baked Appearance}

% Double-anonymized review: author names, affiliations, and acknowledgements
% (including funding) are provided only on the separately uploaded title page.

\begin{abstract}
Geometry-conditioned multi-view diffusion generates views that are combined into textures for existing 3D meshes. The practical value of an inference-time geometry scale depends on the amount actually applied, the conditioning interface and the appearance of the baked asset. We study Timestep-Conditioned Adapter Scaling (TCAS) as an intervention on an existing GeoTex residual path, comparing its historical C3 profile with competitive fixed and linear controls, equal-applied-budget allocations, actual texture baking and native conditioning interfaces. On 300 held-out objects, C3 improves foreground PSNR over Fixed Low by $0.522$ dB, while changing the applied budget; Fixed High and Generic Linear have better means on all four image endpoints. The recorded primary contrast for the retrospective matched-budget timing addendum, Middle-high minus the average of Early-high and Late-high, is not detected ($-0.0339$ dB, 95\% interval $[-0.1172,+0.0514]$, $p=0.4324$). Secondary pairwise comparisons place C3 below Fixed Mean and Late-high by $0.133$ and $0.564$ dB in foreground PSNR. In a ten-object GeoTex bake pilot, Late-high reduces unseen-view foreground LPIPS by $0.00265$, but a separate twenty-object GSO study does not detect this benefit, and substantial appearance errors remain. Native MV-Adapter studies likewise do not detect the specified middle-window advantage. These task-specific findings show that control allocation affects measured outcomes while its practical value depends on protocol and endpoint; they do not establish a universally preferred schedule or faithful material reproduction.
\end{abstract}

\begin{keyword}
\KWD
multi-view diffusion \sep 3D texturing \sep geometry adapter \sep residual scaling \sep timestep conditioning
\end{keyword}

\end{frontmatter}

\raggedbottom
\linenumbers

\section{Introduction}
\label{sec:intro}

Multi-view diffusion can generate a set of images from a reference appearance and known 3D geometry, after which the images are projected or baked onto a mesh. The target is a usable textured asset: generated views should follow the surface while retaining recognizable colors and local appearance when rendered from new cameras. Geometry-adapter strength is an accessible inference-time control within this task.

Evaluating this control requires accounting for how it is implemented and where quality is measured. A temporal profile compared with a weaker constant can also increase total conditioning strength, and wrapper caps can change the requested profile before injection. Generated-view agreement does not by itself measure reference-color fidelity or the appearance of a baked texture, where projection and coverage introduce additional effects. These properties make the choice of geometry control a task-level question that calls for competitive comparisons and asset-level evaluation.

We examine the historical Timestep-Conditioned Adapter Scaling (TCAS) profile C3 on an existing GeoTex-Adapter residual interface. C3 applies a low--high--low request over sampling progress and is evaluated without another schedule search. Scheduling conditioning across time and layers, including geometric conditioning, has prior precedents such as Scheduled Style Injection (SSI) \cite{kulkarni2026ssi}; TCAS is therefore treated as the study intervention, not as a new general scheduling principle. Our focus is the task evidence: what changes under strong and equal-budget controls, whether generated-view differences remain useful after baking, and how findings depend on the conditioning interface.

We address three questions. First, does C3's image-space gain persist against strong controls, and what does a matched applied-scale budget establish about temporal placement? Second, how much of a generated-view effect remains in unseen renders of actual UV-textured meshes, alongside reference appearance and a matched ground-truth bake? Third, do related allocation responses appear under different native conditioning interfaces?

Across 300 held-out objects, C3 improves foreground PSNR over Fixed Low, but the comparison also raises the deep and middle applied-scale sums. Fixed High and Generic Linear have better means on the four reported image endpoints. The original matched-budget primary comparison does not detect a middle-window advantage; secondary pairwise results favor Fixed Mean and Late-high over C3 on several endpoints. A small Late-high LPIPS improvement appears in a ten-object compatible GeoTex bake pilot and is not detected in a separate GSO bake study. Native MV-Adapter studies also do not detect the specified middle-window advantage. A separate historical-request retraining is retained as endpoint-sensitivity evidence with its input-substitution limitation, rather than as confirmation of the main cohort.

The contributions are threefold:

(1) \textbf{A budget-aware characterization of geometry control.} We distinguish requested scales, wrapper-capped applied scales and state-dependent residual corrections. Competitive baselines and matched applied-scale sums show how strength and temporal placement affect generated-view metrics, while the specified primary contrast does not detect the proposed middle-window advantage.

(2) \textbf{An asset-level assessment of generated-view gains.} Actual native-UV meshes, matched-resolution ground-truth self-bakes and eleven unseen cameras quantify a small local gain together with persistent appearance errors and a separate-cohort non-detection. Reference-color agreement, incoming-view consistency and texture variation remain distinct from baked appearance and material fidelity.

(3) \textbf{An evaluation of interface and protocol boundaries.} Native MV-Adapter and retained MVDiffusion studies test related allocation questions under different conditioning operators. Their endpoint-dependent responses define limits to transfer; cohorts are not pooled and identical scale labels are not treated as a common residual dose.

The practical importance is that generated RGB views are intermediate products: a view-level metric change can reflect altered geometry control without improving the textured asset, and a requested scale can be clipped before it reaches the model. For graphics users, the relevant decision is whether changing an inference-time control improves a complete asset under the deployed conditioning interface. This paper supplies bounded, reproducible evidence for that decision across competitive controls, equal-applied-budget timing comparisons and actual UV baking, without proposing a generally optimal schedule.

\section{Related Work}
\label{sec:related}

\subsection{3D Texture Generation and Multi-view Diffusion}
\label{subsec:3dtex}

Single-image 3D generation is often decomposed into geometry generation and texture generation. Single-image-to-3D, multi-view diffusion, and sparse-view reconstruction methods have improved structural stability and view synthesis \cite{liu2023zero1to3,shi2023zero123pp,liu2023syncdreamer,long2023wonder3d,shi2024mvdream,wang2023imagedream,xu2024instantmesh,ke2024era3d}. Texture synthesis remains a distinct challenge for asset realism and usability \cite{perla2026survey}. Many systems generate RGB views with 2D diffusion and project or bake them onto a surface \cite{chen2023text2tex,richardson2023texture,poole2023dreamfusion,chen2023fantasia3d}. Recent work also explores exemplar-based texture synthesis and diffusion-based procedural materials \cite{guzman2026adaptive,lv2025interpretable}.

Paint3D, MVPaint, Hunyuan3D 2.0, and MVPainter use geometry-conditioned multi-view generation in 3D texturing pipelines \cite{zeng2024paint3d,cheng2024mvpaint,lai2025hunyuan3d,shao2025mvpainter}. MV-Adapter, MVDiffusion, and other synchronized multi-view methods address consistency through different adapter and denoising designs \cite{huang2024mvadapter,tang2023mvdiffusion,liu2024synchronized,xu2026wondertex}. These systems establish the graphics context for this work. The present evidence does not compare final baked-texture quality across those systems.

\subsection{Adapter-Augmented Diffusion Models and Geometric Control}
\label{subsec:adapter}

ControlNet and T2I-Adapter provide learned spatial-conditioning branches or feature injections, whose contribution can be weighted by an inference-time condition-strength scale \cite{mou2023t2iadapter,zhang2023controlnet}. They establish conditioning strength as a practical control, but their injected tensors and scaling sites do not share a numerical dose with GeoTex. The studied GeoTex-Adapter path already exposes an inference-time residual scale and evaluates fixed requests. TCAS reuses that control point as a three-stage request over sampling progress, while recording how wrapper-group caps change the applied values and how adapter activations determine the resulting corrections. It adds no learned branch or schedule selector. IP-Adapter introduces image prompts through decoupled attention \cite{ye2023ipadapter}; LoRA and related methods use scaled low-rank updates \cite{hu2022lora,liu2024dora,zhang2023adalora}.

\subsection{Diffusion Guidance Scheduling and Timestep Control}
\label{subsec:guidance}

Diffusion denoising is often discussed in terms of coarse-to-fine evolution \cite{ho2020ddpm,song2021ddim}. Classifier-free guidance combines conditional and unconditional score predictions \cite{ho2022cfg}; limiting it to an interval studies quality and diversity \cite{kynkanniemi2024limited}. Adaptive Guidance instead schedules conditional/unconditional evaluations for inference efficiency \cite{castillo2023adaptive}. These score-space and computation schedules differ from the geometry-residual intervention tested here.

Earlier official ControlNet code already exposes layer-dependent control weights, and Diffusers v0.21.0 documents temporal conditioning start/end gates \cite{controlnet2023layercontrol,diffusers2023controlnet}. SSI is the closest combined layer/time scheduling precedent. It varies StyleID's query-mixing parameter \cite{chung2024styleid} across decoder layers and denoising time, tests schedule direction and shape, and also schedules a separate geometry-conditioned ControlNet branch \cite{kulkarni2026ssi}. Thus, layer/time scheduling and geometry-conditioning schedules are established ideas. TCAS acts on the residual correction produced by GeoTex-Adapter from normal, depth, and foreground-mask conditions, with group-specific wrapper caps. In this 3D texturing setting, we measure how that intervention relates to generated-view agreement, reference-color errors, and baked appearance. The task, interface and endpoints differ from SSI, motivating a task-specific evaluation rather than a claim of a new scheduling principle.

\subsection{Multi-view Adapter Interfaces}
\label{subsec:mvinterfaces}

MV-Adapter uses a condition guider for camera or geometry maps together with decoupled spatial, multi-view, and optional image cross-attention layers \cite{huang2024mvadapter}. MVDiffusion uses correspondence-aware attention modules to connect pixels across views \cite{tang2023mvdiffusion}. These interfaces differ from the GeoTex-Adapter residual branch studied here. We therefore report each auxiliary intervention within its own tested implementation and cohort, without pooling metrics or interpreting it as an architecture-matched replication of TCAS. The native temporal comparisons are summarized in Tables~\ref{tab:new-mva} and \ref{tab:interface-effects}; historical panels remain in Supplementary Section~S3.

\section{Control Formulation and Study Design}
\label{sec:method}

\subsection{Task Definition and Base Pipeline}
\label{subsec:task}

We study inference-time geometry strength in an MVPainter-style multi-view texturing framework \cite{shao2025mvpainter}. Figure~\ref{fig:overview} shows TCAS on the GeoTex-Adapter residual path: the requested scale is limited by the wrapper's group-specific cap before residual injection into frozen UNet features. Given a reference image and 3D geometry, the pipeline generates six RGB target views. Fresh C evaluates these generated views; selected GeoTex and GSO bake cohorts project views onto native UV meshes and assess unlit renders from eleven unseen cameras. The base UNet, VAE and pretrained adapter weights remain frozen.

\begin{figure*}[t]
\centering
\includegraphics[width=\textwidth]{figures/fig1_tcas_control.pdf}
\caption{Task overview of TCAS and asset-level evaluation. The reference, normal, depth, mask and six C3 output views are traceable to the same Building~03 object UID; the source asset is by Shalmon (CC BY~4.0), and the output panel is an example from the recorded C3 confirmation run, not a quality summary. C3 requests $(1.25,2.50,1.25)$ over steps 0--16, 17--32 and 33--49; wrapper-group caps act before residual injection into frozen UNet features. Native-UV baking and unlit renders from eleven unseen cameras summarize selected, separate evaluation cohorts. Baking is an evaluation path, not an added model module.}
\label{fig:overview}
\end{figure*}

During multi-view diffusion, the 3D geometry is rendered into normal maps, depth maps, and foreground masks, denoted as:
\begin{equation}
\mathbf{G} = \{G_\text{normal}, G_\text{depth}, G_\text{mask}\}.
\end{equation}

The geometric adapter maps these conditions to a residual correction term and injects it into the intermediate hidden states of the UNet. For inference step $j$ with normalized execution progress $p_j$, the adapter injection can be written as:
\begin{equation}
\mathbf{h}'_{j,\ell} = \mathbf{h}_{j,\ell} + s_{\mathrm{app}}(p_j,\ell)\, A^\ell_\phi(\mathbf{h}_{j,\ell}, \mathbf{G}),
\label{eq:injection}
\end{equation}
Here $\ell$ indexes an injected layer group. The \emph{requested scale} $s_{\mathrm{req}}(p_j,\ell)$ is the value assigned by the profile; the \emph{applied scale} is the value used by the wrapper,
\begin{equation}
s_{\mathrm{app}}(p_j,\ell)=\min\!\bigl(s_{\mathrm{req}}(p_j,\ell),c_\ell\bigr),
\end{equation}
when the native cap $c_\ell$ is active, and otherwise equals the request. The resulting correction is $\mathbf{r}_{j,\ell}=s_{\mathrm{app}}(p_j,\ell)A^\ell_\phi(\mathbf{h}_{j,\ell},\mathbf{G})$; its \emph{post-scale residual magnitude} is $m_{j,\ell}=\|\mathbf{r}_{j,\ell}\|_2$. Thus neither a requested nor an applied scale label alone determines the correction magnitude.

In the Fresh C wrapper, the three cap groups map the GeoTex residual wrappers on the UNet ResNet blocks as follows: deep comprises the UNet middle block and $\texttt{up\_0}$, middle is $\texttt{up\_1}$, and shallow is $\texttt{up\_2}$. These names identify the implementation's grouped injection points; they do not assert universal semantic roles for denoising stages. The corresponding Fresh C caps are 3.0, 3.5, and 0.8. C3's shallow request of 1.25 is therefore applied as 0.8, while its deep and middle requests pass unchanged. GFL and GFH also clip at the shallow cap. These values describe the evaluated Fresh C wrapper; row-level cap semantics for the historical 24-object probe are not recoverable. TCAS scales a geometry residual added to feature activations; it does not reweight conditional and unconditional score predictions as classifier-free guidance does.

\subsection{Image-Space Metrics and Texture-Variation Diagnostics}
\label{subsec:metrics}

The endpoints describe different outcomes. \emph{Image-space agreement} refers to target-view PSNR, LPIPS, and SSIM measurements on the named image regions; it is not a measure of final 3D material quality. \emph{Reference-color agreement} refers to color differences such as CIEDE2000 computed against the corresponding reference. \emph{Texture variation} refers to statistics such as RGB standard deviation, Laplacian variance, gradient magnitude, and high-frequency energy; these may rise with detail or artifacts and are not fidelity measures. \emph{Subjective preference} refers only to an earlier forced-choice study on its original stimuli. \emph{Baked 3D appearance} refers to rendered mesh outputs after texture projection or baking, while \emph{PBR/material fidelity} additionally concerns appearance under the material and illumination model. These endpoints are not interchangeable. FG-SSIM, Edge-SSIM, and Crop-SSIM quantify structural similarity on the indicated image regions; PSNR measures pixel error; LPIPS compares learned image features \cite{wang2004ssim,zhang2018lpips}.

Laplacian variance, RGB standard deviation, and gradient magnitude can describe local high-frequency or color variation. For example, a foreground Laplacian statistic may be written as
\begin{equation}
M_{\mathrm{var}}=\operatorname{Var}\!\left(\operatorname{Lap}(I_{\mathrm{out}}\odot G_{\mathrm{mask}})\right),
\end{equation}
where $G_{\mathrm{mask}}$ is the foreground mask. A larger value is not necessarily better: detail, noise, repeated patterns, or artifacts can all increase local variation. We therefore do not interpret these proxies as texture fidelity or use them alone to claim that a schedule better preserves reference appearance.

\subsection{TCAS Control Semantics and Comparison Design}
\label{subsec:tcas}

TCAS replaces a globally fixed requested adapter scale with a piecewise profile over normalized sampling progress. The early/middle/late partition is an operational description, not a universal claim about the function of every denoising stage. C3 is the low--high--low profile selected in earlier development and tested here as an empirical intervention, not as an independently derived optimum.

To avoid ambiguity caused by different raw timestep directions in different schedulers, $p_j=j/(T-1)$ denotes execution progress for inference step $j\in\{0,\ldots,T-1\}$. Thus $p=0$ is the first update from the initial noisy latent and $p=1$ is the last denoising update before decoding. For $T=50$, TCAS divides execution into steps 0--16, 17--32, and 33--49 (17/16/17 steps):
\begin{equation}
s_{\mathrm{req}}(p) = \begin{cases} s_e, & 0 \leq p < 1/3 \\ s_m, & 1/3 \leq p < 2/3 \\ s_l, & 2/3 \leq p \leq 1 \end{cases}.
\label{eq:tcas}
\end{equation}

We call this variation across the 50 inference steps the \emph{temporal allocation}. A \emph{layer allocation} varies the requested scale across the three wrapper groups at a fixed step. For a schedule with applied values $s_{\mathrm{app}}(j,\ell)$, its per-group applied-scale sum is $B_\ell=\sum_j s_{\mathrm{app}}(j,\ell)$; the unweighted sum $B_{\mathrm{nom}}=\sum_\ell B_\ell$ is a \emph{nominal applied-scale budget}. It does not weight scales by adapter activations. We log post-scale correction summaries, including wrapper-level $L_2$ norm and mean absolute entry. The three wrapper groups operate on feature maps with different spatial dimensions, and the traces do not retain channel counts; raw $L_2$ values are therefore not used for cross-layer dose comparisons. The per-entry mean absolute correction reported in Supplementary Section~S1 is a dimension-normalized descriptive summary, not an equal-dose intervention. Correction magnitudes also depend on the current hidden state and earlier denoising updates, so equal nominal budgets do not imply equal realized corrections.

The selected C3 configuration used in this paper is:
\begin{equation}
(s_e, s_m, s_l) = (1.25, 2.50, 1.25).
\end{equation}

The low--high--low profile is a concrete inference-time intervention selected during earlier development. Its stage interpretation is a task-specific hypothesis: the tested early, middle, and late allocations may produce different responses. Fresh C evaluates the fixed profile against matched controls but does not establish a universal function for any stage. TCAS adds no trained parameters and does not modify base-model or adapter weights; it replaces the requested scale during inference.

\paragraph{What C3 changes and what its comparisons identify}
C3 requests $(1.25,2.50,1.25)$ for each of the 50 steps across the three layer groups before native caps. In Fresh C, the applied deep and middle profiles retain those values; the shallow profile is clipped to 0.8 throughout. GFL requests 1.25 at every step and group, so its deep and middle values are 1.25 throughout and its shallow value is also clipped to 0.8. Their per-group applied-scale sums are therefore 82.5 versus 62.5 in deep and middle, and 40 for both in shallow. The C3--GFL contrast changes the middle-third deep/middle scale magnitude as well as its temporal allocation; it is not a timing-only comparison.

To identify a temporal-placement contrast, controls hold the per-layer applied values and nominal per-group sums fixed while moving the same values between time windows. A scale-strength contrast instead changes amplitude while preserving layer and temporal allocation. Matching requested values or applied-scale sums does not match the evolving residual corrections; raw $L_2$ values are also dimension-sensitive across wrapper groups and cannot establish cross-layer dose matching. The 300-object allocation comparison matches applied sums in each layer group and therefore estimates an allocation effect at a fixed scalar budget, not a residual-dose-independent mechanism. The historical-request sensitivity panel is reported separately.

\paragraph{Development origin and scope of C3}
C3 $=(1.25,2.50,1.25)$ was selected in a 24-object development probe containing 17 uniform, layer-wise, and time-conditioned variants. The available records do not preserve all variant-level outputs needed to reproduce that selection. Fresh C therefore treats C3 as a fixed operating profile and compares it with named controls without claiming optimality or uniqueness and without further schedule search.

\subsection{Geometry Adapter and Inference Interface}
\label{subsec:adapterimpl}

The pretrained GeoTex-Adapter receives a five-channel input comprising a three-channel normal map, a one-channel depth map, and a one-channel foreground mask. Its multi-scale encoder produces residual features that are injected into the frozen UNet. TCAS acts on the requested multiplier along denoising progress; it does not retrain the adapter, the UNet, or the VAE and introduces no learned parameters.

Visual comparisons use recorded outputs with shared object framing and crops. They illustrate task behavior and appearance limitations; quantitative conclusions use the stated object-level metrics.

\subsection{Study Cohorts, Sampling and Statistical Units}
\label{subsec:setup}

Table~\ref{tab:study-design} summarizes the cohorts and protocols. The principal evaluation cohort, denoted Fresh C, contains 300 object identities disjoint from the historical schedule-selection/evaluation pool, the recorded 1,118-object training list and the earlier clean-v2 evaluation list. Its primary four-condition comparison uses one frozen retrained GeoTex checkpoint, shared inputs and paired object seeds. The historical main checkpoint is unavailable, so Fresh C is a controlled retraining study rather than a direct checkpoint replication. A separate requested-cohort retraining is reported as protocol sensitivity.

\begin{table*}[t]
\centering\scriptsize\setlength{\tabcolsep}{3pt}
\caption{Studies, implemented protocols and evidence scope. Rows are studies rather than disjoint populations: the bake pilots and some interface tests reuse image-study inputs. Object is the statistical unit for validated cohorts; views and seeds are nested within object. The historical-request technical panel instead describes requested slots with one unknown replacement and does not establish independent-object inference.}
\label{tab:study-design}
\begin{tabular}{p{0.16\textwidth}p{0.23\textwidth}p{0.20\textwidth}p{0.32\textwidth}}
\toprule Study / objects & Interface and sampling & Endpoint & Selection and scope \\
\midrule
Fresh C / 300 & Capped GeoTex residual; 256 pixels; 50 Euler steps; conditional-only prediction; cached image features & Six target-view foreground endpoints & UID-disjoint held-out cohort; original primary C3--GFL frozen; linear and equal-budget comparisons are same-cohort addenda \\
GeoTex bake / 10 & Same controlled generator; six 256-pixel bake cameras; native UVs & Eleven unseen unlit views & Input/compatibility-selected subset; primary Late--Mean specified before bake outcomes; two exploratory native controls \\
MV-Adapter / 50 & Native condition-guider features; 512 pixels; 50 steps; CFG 3 & Six native-camera target views & Fresh C inputs not previously viewed through this interface; input-only selection; thirteen legal white targets retained \\
GSO / 100 & Native MV-Adapter; 512 pixels; 50 steps; CFG 3; three seeds & All-six PSNR primary; five nonreference LPIPS secondary & Source-ordered convenience cohort; four profiles had prior baseline exposure; external sensitivity \\
GSO bake / 20 & Same fixed GSO inputs; native UVs; three seeds & Eleven unseen unlit views & Revised GT-only calibration after fifteen original outcomes; full twenty retained; no GT-error subtraction \\
Reference-protocol / 50 & New frozen GeoTex checkpoint; native reference conditions; 512 pixels; 50 steps; CFG 2 & Five nonreference views & Supporting protocol/checkpoint sensitivity on previously studied inputs; not another independent backbone \\
Separate requested cohort / 276 slots & GeoTex retraining; 256-pixel tiles; 50 Euler steps & Full-frame and foreground six-view composites & Previously exposed list; one unrecorded input replacement; endpoint sensitivity only \\
MVDiffusion / 75 & Partial correspondence-block control; depth concatenation unscaled & Recorded six-view endpoints & Retrospective third-interface panel; narrow native-mirror identity check; not exact GeoTex operator transfer \\
\bottomrule
\end{tabular}
\end{table*}

Within each comparison, object inputs, targets, initialization and endpoint code are shared. Fresh C uses seed $42+\texttt{object\_idx}$, held constant across conditions for that object. Its capped residual groups have deep/middle/shallow limits of 3.0/3.5/0.8. MV-Adapter uses a different feature injection operator and numerical scale range; its native target cameras and actual default-path identity are checked separately. Scale labels are not common physical doses across interfaces.

Fresh C's controlled sampler differs from the complete base pipeline, which uses online image features, classifier-free guidance and 512-pixel views; it has no input- and seed-matched complete-default condition. A two-object default-path identity check is reported separately in Supplementary Section~S6 and is not a quality comparison. Constant 1.0 under the GeoTex caps also differs from the unset-scale residual branch, so neither it nor No Adapter is called the native pipeline. We name controls by their executed configuration; source records and sampling/default-path diagnostics are in Supplementary Sections~S1 and S6.

FG-PSNR, LPIPS, SSIM and Edge-SSIM measure agreement on the named image regions. Baked evaluation renders actual exported meshes; unlit base color does not establish PBR or material fidelity. Fresh C image endpoints are evaluated once on each frozen $3\times2$ six-view composite. Native-interface and baked endpoints average their recorded per-view scores within object, with GSO sampling seeds also averaged within object. Paired intervals resample objects and are pointwise 95\% intervals. Preserved primary, secondary and exploratory families use their stated Holm adjustments; the distinct families do not imply one overall experiment-wide error guarantee. Non-detection is not equivalence or non-inferiority. Complete comparison rules are supplied in Supplementary Sections~S2--S5.

\section{Experimental Results}
\label{sec:experiments}

\subsection{Competitive Controls and the Historical C3 Profile}
\label{subsec:largescale}

The internally specified Fresh C primary is C3 minus Fixed Low (GFL) FG-PSNR. The paired gain is $+0.52$ dB (95\% interval $[+0.40,+0.64]$), with 206/300 objects favorable. All conditions use the same frozen retrained checkpoint and object inputs. This gain combines increased strength and temporal allocation: C3's deep/middle applied sums are 82.5 each, versus 62.5 for GFL, while both have shallow sum 40.

\begin{table*}[t]
\centering\scriptsize\setlength{\tabcolsep}{4pt}
\caption{Fresh C competitive-control means on the same 300 objects. Early/Middle/Late-high and Fixed Mean share exactly the same per-group applied sums. The equal-budget study is a post-inspection addendum, not a second holdout. Each budget entry for deep/middle applies separately to both groups. $^\dagger$Generic Linear is a competitive same-cohort addendum, not part of the matched-budget family; its different group budgets are shown explicitly.}
\label{tab:freshc-absolute}\label{tab:matched-budget300}
\begin{tabular}{lrrrrrr}
\toprule Control & Deep/middle sum & Shallow sum & FG-PSNR$\uparrow$ & FG-LPIPS$\downarrow$ & FG-SSIM$\uparrow$ & Edge-SSIM$\uparrow$ \\
\midrule
\input{tables/matched_controls014.tex}
\bottomrule
\end{tabular}
\end{table*}

Competitive controls change the recommendation. Fixed High has better means than C3 on all four endpoints: C3 minus Fixed High is $-0.748$ dB FG-PSNR (pointwise interval $[-0.885,-0.609]$), with FG-SSIM difference $-0.04553$ ($[-0.04937,-0.04167]$). Generic Linear ramps deep/middle requests from 1.25 to 2.50 and shallow requests from 0.50 to 0.75 over the fifty steps. Its applied sums are (93.75,93.75,31.25), so it changes layer allocation and total budget as well as timing. It also has better means on all four; C3 minus Generic Linear is $-0.593$ dB ($[-0.837,-0.351]$). Its positive median of $+0.074$ dB and 155/300 favorable objects show why mean and object-level distributions must both be considered. Complete original paired endpoints and their comparison families are in Supplementary Section~S2.

\subsection{Temporal Allocation at Matched Applied-Scale Budgets}
\label{subsec:matched-budget300}
Early-, middle-, and late-high move the same sixteen high steps to updates 0--15, 17--32, and 34--49. Deep/middle low and high are 1.25 and 2.50; shallow stays at its 0.8 cap. Fixed Mean uses 1.65 before the cap. Each allocation therefore has per-group sums $(82.5,82.5,40)$. The completed six-condition matrix contains 1,800 paired rows with common-input and seed checks.

The protocol-defined primary contrast for this retrospective timing addendum is Middle-high minus the objectwise average of Early-high and Late-high FG-PSNR. Its mean is $-0.0339$ dB (10,000-resample pointwise object-bootstrap interval $[-0.1172,+0.0514]$; two-sided 20,000-draw sign-flip $p=0.4324$), with 119/300 objects positive. The median is $-0.2296$ dB. This result does not support a middle-window advantage and does not establish equivalence. The archived protocol records that this contrast was specified before the new allocation outputs; earlier Fresh C outcomes were already known, so it is not an independent prospective confirmation. A later all-pair recomputation reports Middle-high versus Fixed Applied Mean under a separate descriptive family; both analysis records and their chronology are distinguished in Supplementary Section~S2.

C3 minus Fixed Mean is $-0.133$ dB (pointwise interval $[-0.170,-0.096]$); C3 minus Late-high is $-0.564$ dB ($[-0.634,-0.494]$). C3 also trails both on LPIPS, FG-SSIM and Edge-SSIM, while improving over Early-high by approximately 0.50 dB in FG-PSNR. The objectwise winner counts and post-inspection four-endpoint dominance summaries are descriptive and are reported in Supplementary Section~S2; they do not imply statistical separation on each object.

The equal-budget comparison identifies an allocation total effect. Evolving residual amplitudes remain trajectory-dependent, preventing a residual-dose-independent mechanism interpretation. Late-high is not universally preferred: Fixed Mean has higher FG-SSIM, and Generic Linear is a strong unmatched-budget practical control. The historical LLH layer/time addendum and complete object distributions remain in Supplementary Section~S2.

\begin{figure*}[!t]
\centering
\includegraphics[width=\textwidth]{figures/fig5_freshc_context_roi.pdf}
\caption{Fresh C view 0: (a) Octa\`edre, a fixed-scale endpoint trade-off; (b) Cube 01, a Generic Linear example. Each column shows the complete object in a common frame and the same pixel ROI enlarged below. Object-level C3 minus Generic Linear in (b) is $-3.491$~dB FG-PSNR and $+0.08172$ FG-LPIPS, favoring Linear on both. The added patterns differ visibly from the low-variation target. These outcome-selected examples illustrate behavior, not prevalence or material fidelity. Composite object metrics use all six views, not only the displayed crop. Object credits are in Supplementary Section~S9.}
\label{fig:c3compare}
\end{figure*}


\subsection{Appearance and Color Failure Analysis}
\label{subsec:appearance-color}

The selected color diagnosis contains 26 failure-enriched objects: six earlier figure cases and 20 Fresh C panels. It is not a random sample, a new holdout, or a population estimate. No Adapter, GFL, and LLH were compared on the same inputs, checkpoint, and runner with six target views per object (78 object-condition records). The following paired object-bootstrap intervals are descriptive (20,000 resamples; seed 61008), and no confirmatory multiplicity family was defined.

\begin{figure*}[t]
\centering
\includegraphics[width=\textwidth]{figures/fig6_color_diagnostic_freshc_context_roi.pdf}
\caption{Fresh C view 0, E5-5-2: target, Fixed Low, C3 and layer LLH share a complete-object frame and an identical pixel crop. The generated brown pattern does not reproduce the gray target; changing its variation does not establish color fidelity. This previously viewed example came from an existing C3--GFL LPIPS stratum and was not selected for color. It does not reproduce the legacy purple-shift cases. Object credits are supplied in Supplementary Section~S9.}
\label{fig:color-failure}
\end{figure*}

\begin{table}[t]
\centering
\scriptsize
\setlength{\tabcolsep}{3pt}
\caption{Selected paired results from the failure-enriched 26-object color diagnosis. Negative CIEDE2000 differences indicate lower reference-color error for the first condition. The set is descriptive and does not estimate population prevalence.}
\label{tab:color-summary}
\begin{tabular}{p{2.3cm}p{3.3cm}p{1.3cm}}
\toprule
Contrast & Mean CIEDE2000 difference [95\% CI] & Objects \\
\midrule
GFL -- No Adapter & $-19.654$ [$-26.370,-13.248$] & 24/26 lower \\
LLH -- GFL & $+3.583$ [$+0.462,+6.712$] & -- \\
\bottomrule
\end{tabular}
\end{table}

GFL lowered reference-relative CIEDE2000 in 24/26 objects, while visible pink or purple shifts remained in two previously documented examples under No Adapter, GFL, and LLH. Their source panels cannot be matched to the preserved object-level outputs; the Fresh C example in Figure~\ref{fig:color-failure} is a separate object. The shifts are present in decoded RGB before UV projection or baking. In one case, a ground-truth VAE encode--decode isolate had CIEDE2000 4.086 and $\Delta a^*\approx-1.35$, compared with 22.032 and $+1.977$ for GFL. This weakens ordinary VAE round-trip error as the sole explanation for that example, without identifying the complete cause.

LLH--GFL increased mean CIEDE2000 while reducing foreground gradient magnitude, high-frequency energy and color entropy. Lower variation therefore did not imply closer reference appearance. The denoising/runtime isolates and three failed correction candidates remain in Supplementary Section~S6; they localize observations but do not validate a cause or repair.

Earlier bake results and calibration checks are reported separately in Supplementary Section~S5 and are not pooled with the ten-object cohort. The GeoTex residual branch with an unset scale differs from an explicit Constant 1.0 request under the wrapper caps; we label the latter by its executed setting.

\paragraph{Full-object evidence and its selection.}
The strong-fixed and Generic Linear examples in Figure~\ref{fig:c3compare} and the color failure in Figure~\ref{fig:color-failure} use common framing and exact shared pixel crops. Favorable/adverse C3--GFL examples are retained in Supplementary Section~S9. Outcome-selected displays illustrate behavior, not prevalence. The historical texture audit does not preserve a traceable denominator or validated texture-loss labels; we therefore do not infer color-failure prevalence or texture loss from it. Variation statistics alone cannot establish fidelity.


\subsection{Actual Baked Appearance: Local Benefit and External Limits}
\label{subsec:new-bake}
We selected ten previously unbaked Fresh C objects in UID order using input availability, native UV/material compatibility, triangle count and camera-mask alignment. Only ten of the 300 screened inputs passed these criteria before the cutoff. Their generated-view outputs were known, but their baked outcomes had not been inspected during selection. This compatibility-selected pilot is not a representative sample; five earlier bake objects are reported separately.

Seven conditions share each object's source mesh, reference, seed, six projection cameras, masks, 256-pixel source resolution, atlas resolution, hole filling and eleven unseen evaluation cameras. Original geometry and UVs are retained; unseen cameras do not enter the bake. GT-256 uses the same source resolution, masks and bake parameters, with the same raw atlas coverage as each generated condition. Camera and sampling calibrations are documented in Supplementary Section~S5. Evaluation uses unlit base color and does not measure PBR fidelity.

\begin{table}[ht]
\centering\scriptsize
\caption{Ten-object unlit unseen-view bake pilot. Eleven views are averaged per object. Higher PSNR/SSIM and lower LPIPS/CIEDE2000 are favorable. Constant 1.0 is an explicitly capped request, not the unset-scale pipeline. All seven planned schedule conditions and all ten eligible objects are retained; the two exploratory controls marked $\dagger$ are auxiliary; GT-256 is a separate matched reference bake, not a theoretical bound.}
\label{tab:new-bake}
\begin{tabular}{lrrrr}
\toprule Condition & FG-PSNR & FG-LPIPS & FG-SSIM & CIEDE2000 \\
\midrule
GT-256 reference bake & 22.38 & 0.0590 & 0.7683 & 4.44 \\
\midrule
Native residual$^\dagger$ & 11.13 & 0.1372 & 0.3376 & 26.47 \\
No Adapter$^\dagger$ & 5.92 & 0.1600 & 0.3829 & 40.92 \\
Constant 1.0 (capped) & 11.93 & 0.1354 & 0.3504 & 24.18 \\
Fixed Low & 12.48 & 0.1317 & 0.3835 & 22.71 \\
Fixed High & 12.29 & 0.1305 & 0.4587 & 23.27 \\
Fixed Mean & 12.62 & 0.1303 & 0.4193 & 22.32 \\
C3 / Middle-high & 12.48 & 0.1302 & 0.4059 & 22.16 \\
Late-high & 12.92 & 0.1277 & 0.4259 & 21.65 \\
Generic Linear & 11.60 & 0.1276 & 0.4458 & 24.71 \\
\bottomrule\end{tabular}
\end{table}

The internally predeclared primary contrast was Late-high minus Fixed Mean on unseen-view FG-LPIPS, proposed from the earlier pilot and tested on the ten newly selected bake objects. Its mean is $-0.00265$ (95\% object-bootstrap CI $[-0.00414,-0.00137]$; exact paired sign-flip $p=0.00195$), with all ten objects favorable. This is approximately a 2.0\% reduction relative to Fixed Mean's mean LPIPS, and identifies a small allocation benefit within this sampling protocol. Generic Linear has a slightly lower mean LPIPS but lower PSNR; Late-high does not dominate every endpoint or control. The matched GT-256 reference bake has mean LPIPS 0.0590, versus 0.1303 for Fixed Mean and 0.1277 for Late-high. A post hoc description therefore puts the absolute gain in context: it closes only 3.7\% of the mean Fixed Mean--GT reference gap (pointwise object-bootstrap CI [2.3\%,5.3\%]). This ratio of paired means is not a theoretical ceiling, a new primary test or evidence of acceptable material fidelity. Leaving out each object in turn keeps the mean Late--Mean difference between $-0.00294$ and $-0.00208$, but all substantial reference-appearance discrepancies remain. Supplementary Section~S5 supplies the paired values and calculation.

C3--Fixed High has mean PSNR $+0.19$~dB (95\% CI $[-0.65,+0.87]$) and LPIPS $-0.00036$ ($[-0.00396,+0.00336]$); neither separates under the four-comparison secondary Holm family (both adjusted $p=1$). Its SSIM is lower by 0.05274 ($[-0.07459,-0.03170]$; Holm $p=0.00781$). Thus the old five-object C3 PSNR advantage does not repeat stably on this new pilot. We retain C3 as the historical candidate and do not rename Late-high as C3.

Across the seven conditions, mean within-object Spearman correlations between generated-view and baked unseen-view rankings are 0.846 for PSNR, 0.804 for LPIPS, 0.746 for SSIM and 0.907 for Edge-SSIM. Broad rankings are positively associated, while the winning condition changes on 7/10 objects for LPIPS (1/10 for PSNR). These are different statements: we do not claim a general reversal of rankings. Full-object Figure~\ref{fig:baked-context} also shows persistent appearance errors, including the Tire's brown output against a gray reference. A lower LPIPS does not establish faithful material reproduction.

\begin{figure*}[t]
\centering
\includegraphics[width=\textwidth]{figures/fig8_baked_unseen_context.pdf}
\caption{Unseen camera 004: (a) LG Remote and (b) Tire. Source and matched GT-256 bake separate projection loss from generated appearance; Fixed Mean and Late-high show the specified allocation comparison. Columns share a full-object frame across conditions. Persistent brown-pattern and material errors remain visible. Complete seven-condition displays for the rights-cleared examples are in Supplementary Section~S5; all ten objects enter the numerical analysis. Credits: mohd.akmal.mahyudin and Koto3D Stephane Vogley, CC Attribution (source links in Supplementary Section~S5).}
\label{fig:baked-context}
\end{figure*}


\paragraph{Separate twenty-object GSO evaluation.}
The GSO/MV-Adapter study completes twelve conditions by three seeds on twenty fixed inputs, plus twenty matched GT self-bakes: 740 unique baked meshes and 8,140 unseen-view records. The objects, source meshes, projection cameras, texture resolution and eleven unseen cameras are shared across controls; unseen cameras do not contribute to baking.

An initial UV-ramp calibration stopped after fifteen objects when a discontinuous diagnostic failed. A subsequent ground-truth-only calibration addressed finite-resolution mixing across UV boundaries and retained the fixed twenty-object cohort, cameras, masks, materials and resolutions. Because this correction followed inspection of initial outcomes, the resulting GSO evaluation is reported as protocol sensitivity. The original failure and remaining UV/coverage limits are detailed in Supplementary Section~S4.

Late-high minus Fixed Mean unseen-view LPIPS is $+0.000084$ (pointwise interval $[-0.000103,+0.000286]$; original comparison-family Holm $p=0.788$). This study does not detect the ten-object pilot's benefit. It changes both input domain and conditioning/projection protocol, so it is neither a direct same-protocol replication nor evidence assigning the difference to backbone alone. Table~\ref{tab:bake-cohorts} presents both application comparisons.

\begin{table*}[t]\centering\scriptsize\setlength{\tabcolsep}{4pt}
\caption{Late-high minus Fixed Mean on unseen-view foreground LPIPS. Lower is favorable. All views/seeds are nested within object. The ten-object specified primary uses its original exact sign-flip test; GSO uses the original Holm16 comparison family. Pointwise intervals are not simultaneous. The two cohorts and protocols are not pooled.}
\label{tab:bake-cohorts}
\begin{tabular}{lrrrrl}
\toprule Study & Objects & Fixed Mean & Late-high & Paired difference [95\% CI] & Test \\
\midrule
GeoTex compatibility pilot & 10 & 0.1303 & 0.1277 & $-0.002652\ [-0.004144,-0.001367]$ & exact $p=0.00195$ \\
GSO revised-calibration study & 20 & 0.10089 & 0.10098 & $+0.000084\ [-0.000103,+0.000286]$ & Holm $p=0.788$ \\
\bottomrule\end{tabular}\end{table*}

Across all twelve GSO conditions, mean within-object image/bake LPIPS-rank correlation is 0.601 (object interval $[0.440,0.733]$). A post hoc sensitivity restricted to the seven original scale/time controls gives 0.396 ($[0.175,0.582]$), with the minimum-mean LPIPS condition changing on 11/20 objects. These are positively associated rankings with changing point-estimate winners, not general rank reversal or statistically resolved winner changes. The GT self-bake has mean LPIPS 0.0602 and PSNR 19.36 dB; average complete-bilinear-footprint coverage is about 78.7\%. All foreground remains scored and GT errors are not subtracted. These losses contextualize the small method effects without defining a theoretical quality bound.

\paragraph{Incoming-view consistency.}
The bake also measures cross-camera color variance at UV texels receiving multiple source views. A constant-black control has zero variance despite its incorrect appearance, showing that this diagnostic measures incoming-view disagreement rather than fidelity. Late-high has a lower mean disagreement than Fixed Mean ($-0.002448$), while Fixed High has still lower disagreement and higher baked LPIPS. Resampling and overlapping UVs also affect this measure; it does not validate seams or PBR. Full controls and contrasts are in Supplementary Section~S5.

\paragraph{Native controls and protocol sensitivity.}
Native residual and No Adapter outputs on the same ten GeoTex objects are exploratory controls. They retain matched geometry and GT coverage, and the repeated Constant 1.0 (capped) outputs are bitwise identical. Their apparent PSNR/LPIPS gains for the schedules do not separate under the fifteen-test exploratory family. The complete data are in Supplementary Section~S5.

A supporting fifty-object comparison uses a newly trained frozen GeoTex checkpoint with 512-pixel supervision, actual positive/negative reference conditions and CFG 2. Late-high minus Fixed Mean has nonreference PSNR $-1.184$ dB (interval $[-1.587,-0.849]$) and LPIPS $+0.00266$ ($[-0.00003,+0.00640]$; adjusted $p=0.080$). C3 has no corrected advantage over Fixed Mean or Fixed High. This previously studied cohort and changed checkpoint/protocol supply a sensitivity boundary, not an independent backbone confirmation. Training and full-v29 strong-baseline results remain in Supplementary Section~S8.

\subsection{Transfer to other conditioning interfaces}
\label{subsec:cross-interface}
MV-Adapter and MVDiffusion supply different conditioning operators and scale semantics from GeoTex. We evaluate each within its native target/camera contract and do not pool their scores or interpret identical scale labels as common residual doses.

The retrospective 75-object MVDiffusion panel scales correspondence-block outputs rather than the GeoTex residual operator. Its layer-wise L--LLH condition is worse than global fixed low in PSNR, structural agreement and reference-color error; LPIPS does not separate. This result and the earlier MV-Adapter panel remain in Supplementary Section~S3 and delimit the scope of cross-interface transfer.

\paragraph{New native MV-Adapter temporal comparison.}
\label{subsec:new-mva}
We selected fifty Fresh C inputs not previously evaluated through MV-Adapter, using input availability and native material/geometry compatibility. Native mesh normalization and six orthographic cameras define both generation and evaluation targets. Three-object identity checks confirmed agreement between the deployed and recorded default paths. An earlier seven-object panel was superseded after a camera-calibration error and is not pooled with this study.

The seven conditions share object inputs, seed, geometry and evaluation code. With 50 steps, low/high are 0.75/1.0; stage-high schedules contain 16 high steps, Fixed Mean is 0.83, and Generic Linear is 0.75--0.91. These five controls have matched per-feature sums within this interface. All 50 objects remain in the primary cohort, including thirteen texture-absent sources with legal constant-white base-color targets. Targets are unlit native-camera base color, not Fresh C's differently normalized reference renders or certified physical albedo.

\begin{table}[ht]\centering\scriptsize
\caption{Native-camera MV-Adapter comparison, fifty input-selected objects. Views are averaged within object. Scale values and endpoints are specific to this interface, and are not numerical doses shared with GeoTex.}\label{tab:new-mva}
\begin{tabular}{lrrrr}
\toprule Condition & FG-PSNR & FG-LPIPS & FG-SSIM & Edge-SSIM \\
\midrule
Native / Fixed High & 12.60 & 0.1280 & 0.5041 & 0.3526 \\
Fixed Low & 12.73 & 0.1270 & 0.5066 & 0.3475 \\
Fixed Mean & 12.68 & 0.1279 & 0.5068 & 0.3509 \\
Early-high & 12.60 & 0.1296 & 0.5041 & 0.3503 \\
Middle-high & 12.71 & 0.1262 & 0.5061 & 0.3470 \\
Late-high & 12.74 & 0.1260 & 0.5075 & 0.3514 \\
Generic Linear & 12.72 & 0.1263 & 0.5072 & 0.3501 \\
\bottomrule\end{tabular}\end{table}

The internally predeclared primary, Middle-high minus Fixed Mean FG-PSNR, is $+0.029$~dB (95\% object-bootstrap CI $[-0.019,+0.081]$; paired Monte Carlo sign-flip $p=0.281$; 29/50 favorable). This does not support a stable middle-window gain or establish equivalence. The remaining 104 pair--metric tests use one Holm family. Middle-high has lower Edge-SSIM than Fixed Mean (adjusted $p=0.00104$); some other placement contrasts separate, including Late-high versus Early-high PSNR. Such endpoint-dependent differences do not demonstrate transferable C3 optimality. Complete contrasts, correction rules and the textured subset are in Supplementary Section~S3. This is a native-interface validation of the scheduling question, not architecture-matched replication of GeoTex's capped residual intervention.


\begin{figure*}[t]
\centering
\includegraphics[width=\textwidth]{figures/fig9_mvadapter_native014.pdf}
\caption{Native MV-Adapter view 0: (a) LG Remote, (b) Tire. Target, Fixed Mean, Middle-high and Late-high use the native camera, object seed and a common complete-object frame. The remote preserves recognizable details while the tire retains an appearance mismatch; neighboring scheduling outputs are visually similar. The fifty-object primary interval includes zero. All seven controls and object credits are provided in Supplementary Sections~S3, S5 and S9.}
\label{fig:native50-visual}
\end{figure*}

\paragraph{External-input native schedule comparison.}
The source-ordered GSO100 convenience cohort uses the same native MV-Adapter control path and three shared seeds. It includes seven scale/time profiles, with four inherited profiles frozen after earlier baseline exposure. The scheduling comparison is an external protocol sensitivity, not an independent prospective confirmation or parameter search.

The specified all-six-view Middle-high minus Fixed Mean primary is $-0.014$ dB (interval $[-0.031,+0.004]$; paired $p=0.122$). A secondary Late-high--Fixed Mean contrast is $+0.045$ dB ($[+0.024,+0.065]$; Holm83 $p=0.00335$), with SSIM $+0.00134$ (adjusted $p=0.000415$). Its five-nonreference-view LPIPS difference is $-0.000086$ ($[-0.000551,+0.000413]$; separate Holm21 $p=1$). Thus a small detected image effect does not establish a perceptual or baked-appearance benefit. Table~\ref{tab:interface-effects} and Figure~\ref{fig:allocation-effects} put these named comparisons in context.

\begin{table*}[t]\centering\scriptsize\setlength{\tabcolsep}{3pt}
\caption{Selected within-study paired effects. P: the specified primary, unadjusted within its study; S: secondary with its original family adjustment (GSO all-six Holm83, nonreference LPIPS Holm21, reference-protocol secondary Holm18). Intervals are pointwise. Native50 and GSO100 use different inputs/seed policies; the reference-protocol study also changes weights and sampling and is not a separate backbone.}
\label{tab:interface-effects}
\begin{tabular}{lllrrcr}
\toprule Study & Contrast & Endpoint & Mean & 95\% CI & Class & $p$ \\
\midrule
\input{tables/interface_effects014.tex}
\bottomrule\end{tabular}\end{table*}

\begin{figure*}[t]\centering
\includegraphics[width=\textwidth]{figures/fig_allocation_effects014.pdf}
\caption{Paired allocation effects from completed object-level records. Error bars are pointwise 95\% object intervals. Marker/color keys distinguish the recorded timing primary, the later all-pair analysis, native-interface results and protocol sensitivity. In panel (a), the first GeoTex row is the protocol-defined Middle-high minus mean(Early-high, Late-high) timing primary; the second is a separate later Middle-high--Fixed Applied Mean contrast. Other rows use Fresh C composite PSNR or native-interface mean-view PSNR; the reference-protocol row excludes its reference view. Panel (b) uses eleven unseen baked views; the GSO20 row is a separate external cohort. Studies are shown separately, without pooling or cross-pipeline ranking.}
\label{fig:allocation-effects}
\end{figure*}

\paragraph{Fixed-seed sensitivity.}
A post-inspection analysis retains all three recorded sampling seeds. GSO100 Late--Mean PSNR averages are positive at each seed ($+0.057,+0.038,+0.039$ dB), but only 27/100 objects improve at all three and 68/100 have mixed or tied directions. Its nonreference LPIPS averages change direction ($-0.000389,-0.000269,+0.000398$). For the GSO20 bake, only 1/20 objects has lower Late--Mean LPIPS at all three seeds; 14/20 have mixed or tied directions. Sign agreement has no practical-effect threshold. These conditional, object-level summaries do not estimate training randomness or revise the original primary tests; full seed intervals and deltas are in Supplementary Section~S4.

Projection-guidance candidates are retained as exploratory negative results in Supplementary Sections~S4 and S8; they do not form a validated alternative method.

\subsection{Endpoint Sensitivity in a Separate Requested Cohort}
\label{subsec:original276}

An auxiliary retraining on a previously exposed 276-request list illustrates how evaluation region can change the conclusion. Its slot-level C3--Fixed High contrast is $+1.699$ dB in full-frame PSNR but $-1.314$ dB in foreground PSNR; FG-LPIPS is higher by 0.01127, FG-SSIM lower by 0.05961, and Edge-SSIM lower by 0.02777. An empty-input failure at requested slot 46 triggered a random replacement whose UID was not recorded. The list therefore does not establish 276 distinct held-out objects, and this panel is protocol/checkpoint sensitivity rather than confirmation of the main cohort or a test of temporal allocation. Condition means, complete paired contrasts and the failed-slot sensitivity are reported in Supplementary Section~S7.


\FloatBarrier
\section{Discussion and Limitations}
\label{sec:limitations}

\paragraph{What the controls establish.}
The Fresh C300 gain over Fixed Low supports the recorded intervention under its frozen protocol, with both applied budget and timing changing. The equal-applied-budget matrix estimates a temporal-allocation effect while retaining trajectory-dependent residual corrections. Its specified Middle--average(Early,Late) primary does not detect the proposed middle-window advantage. Secondary pairwise contrasts favor Fixed Mean and Late-high over C3 on foreground PSNR, with endpoint trade-offs across the remaining measures. The overlap-excluded historical-request retraining shows opposing full-frame and foreground C3--Fixed High directions, while its unknown input substitution limits inference to a protocol sensitivity analysis.

\paragraph{Application value and operating points.}
In the GeoTex image protocol, Fixed High and Generic Linear have stronger endpoint means than C3; equal-budget late-high trades higher PSNR and lower LPIPS against lower FG-SSIM than fixed mean. The ten-object bake's approximately 2\% LPIPS reduction closes only 3.7\% of the gap to the matched GT reference bake, while visible appearance errors remain. The separate GSO20 study does not detect that allocation benefit. Its post hoc calibration, GT self-bake loss and incomplete coverage are disclosed; a small image effect alone does not justify replacing an existing deployment setting.

The evidence suggests a practical evaluation sequence rather than a new selection algorithm: compare native/default and competitive fixed settings inside the same deployed protocol; use matched applied budgets to assess timing; check reference appearance and the actual unseen-view bake; and repeat the named comparison across fixed seeds before claiming repeatable objectwise improvement. The tested GeoTex image setting supports stronger fixed controls than historical C3. The native-interface and GSO bake results provide no general reason to replace a deployment's fixed setting with the middle window. These bounded conclusions do not depend on a universally winning schedule.

Generated-view metrics remain informative: overall baked/image method rankings are positively associated. However, preferred point estimates often change, and the association is weaker when the GSO analysis is restricted to closely competing scale/time settings. Winner changes do not prove significant separation between methods or a general reversal. Actual baked evaluation and full-object reference comparison remain necessary.

\paragraph{Interface and protocol boundaries.}
Native MV-Adapter50 and GSO100 do not detect the primary middle-window advantage. The secondary GSO late-window PSNR change is statistically detectable but small, with no corresponding nonreference LPIPS benefit. The supporting reference-conditioned checkpoint shows a different Late--Mean direction. Historical MVDiffusion supplies a partial-operator boundary. Because weights, inputs, conditioning operators and sampling protocols differ, these studies cannot isolate one backbone or domain as the cause of the varying effects.

\paragraph{Quality, selection and replication limits.}
The bake studies use compatible native-UV assets, and GSO uses a source-ordered convenience sample. They do not establish representative asset quality, physical albedo, PBR fidelity or a validated seam score. Views and seeds are nested within object, and seed repetitions do not estimate training-checkpoint variation. Additional seed, dominance and UV-overlap summaries are descriptive analyses defined after outcomes were observed. Separate test families and pointwise intervals retain their stated roles; non-detection is not equivalence.

Fresh C uses a controlled retraining and sampler, so its estimates do not identify performance for an unauthenticated historical checkpoint or the complete default pipeline. The separate requested-cohort panel has one unknown substituted input and cannot support inference over a verified 276-object cohort. These limits qualify cross-checkpoint replication claims while leaving the named within-protocol paired comparisons intact.

\section{Conclusion}
\label{sec:conclusion}

TCAS profiles produce different generated-view and baked outcomes within the tested implementations. On 300 held-out objects, C3 improves FG-PSNR over Fixed Low, while this contrast also increases the applied deep and middle scales; Fixed High and Generic Linear have better means on all four image endpoints. At matched applied-scale sums, the specified middle-versus-early/late-average primary does not detect a middle-window advantage. Secondary pairwise contrasts favor Fixed Mean and Late-high over C3 on foreground PSNR. In actual baking, Late-high yields a small LPIPS improvement on ten compatible GeoTex objects, while a separate twenty-object GSO study does not detect the same benefit and visible appearance errors remain. Native MV-Adapter studies likewise do not detect the specified middle-window advantage. These task-specific findings connect denoising controls to textured-asset evaluation and do not establish a universally preferred schedule or faithful material reproduction.

\section*{Data availability}

The accompanying research package contains object-level metrics, frozen protocols and configuration records, analysis scripts, and the LaTeX source needed to rebuild the manuscript and Supplement. Recovered FAC source is supplied with a record of its entry points and validation scope; these component checks do not constitute a historical full-model replay. Statistical summaries can be recomputed without model weights, while fresh inference requires the recorded weights, runtime and licensed source assets. A supplementary 5.5-second video shows one matched GSO object across the recorded bake conditions and eleven unseen cameras; it is illustrative and does not replace object-level analysis. Raw generated-image archives, trained checkpoints and participant-level responses are not included. Source assets must be acquired under their recorded licenses. Additional analysis records may be requested from the corresponding author, subject to licensing, privacy and archive constraints.

\section*{Declaration of generative AI and AI-assisted technologies in manuscript preparation}
OpenAI Codex assisted with revision drafting, implementation and reproducible statistical and figure code. The final interface diagram uses reproducible vector elements; research displays come from the recorded runs without image-generation retouching. AI-assisted presentation checks cover rendered figures and pages for legibility; these checks do not certify scientific validity or replace author review. Numerical tables are computed from delivered object-level records. Final scientific review, accuracy and responsibility remain with the authors.

\section*{Declaration of competing interests}

The authors declare that they have no known competing financial interests or personal relationships that could have appeared to influence the work reported in this paper.


\begingroup
\linespread{0.94}\selectfont
\begin{thebibliography}{45}
\bibitem{liu2023zero1to3}
Liu R, Wu R, Van Hoorick B, et al. Zero-1-to-3: Zero-shot one image to 3D object. \textit{Proceedings of the IEEE/CVF International Conference on Computer Vision}, 2023.

\bibitem{shi2023zero123pp}
Shi R, Chen H, Zhang Z, et al. Zero123++: A single image to consistent multi-view diffusion base model. arXiv:2310.15110, 2023.

\bibitem{liu2023syncdreamer}
Liu Y, Lin C, Zeng Z, et al. SyncDreamer: Generating multiview-consistent images from a single-view image. \textit{International Conference on Learning Representations}, 2024.

\bibitem{long2023wonder3d}
Long X, Guo Y C, Lin C, et al. Wonder3D: Single image to 3D using cross-domain diffusion. \textit{Proceedings of the IEEE/CVF Conference on Computer Vision and Pattern Recognition}, 2024.

\bibitem{shi2024mvdream}
Shi Y, Wang P, Ye J, et al. MVDream: Multi-view diffusion for 3D generation. \textit{International Conference on Learning Representations}, 2024.

\bibitem{wang2023imagedream}
Wang P, et al. ImageDream: Image-prompt multi-view diffusion for 3D generation. arXiv:2312.02201, 2023.

\bibitem{xu2024instantmesh}
Xu J, et al. InstantMesh: Efficient 3D mesh generation from a single image with sparse-view large reconstruction models. \textit{Advances in Neural Information Processing Systems}, 2024.

\bibitem{ke2024era3d}
Li P, Liu Y, Long X, et al. Era3D: High-resolution multiview diffusion using efficient row-wise attention. \textit{Advances in Neural Information Processing Systems}, 2024.

\bibitem{chen2023text2tex}
Chen D Z, et al. Text2Tex: Text-driven texture synthesis via diffusion models. \textit{Proceedings of the IEEE/CVF International Conference on Computer Vision}, 2023.

\bibitem{richardson2023texture}
Richardson E, Metzer G, Alaluf Y, et al. TEXTure: Text-guided texturing of 3D shapes. \textit{ACM SIGGRAPH 2023 Conference Proceedings}, 2023.

\bibitem{poole2023dreamfusion}
Poole B, Jain A, Barron J T, Mildenhall B. DreamFusion: Text-to-3D using 2D diffusion. \textit{International Conference on Learning Representations}, 2023.

\bibitem{chen2023fantasia3d}
Chen R, Chen Y, Jiao N, Jia K. Fantasia3D: Disentangling geometry and appearance for high-quality text-to-3D content creation. \textit{Proceedings of the IEEE/CVF International Conference on Computer Vision}, 2023.

\bibitem{shao2025mvpainter}
Shao M, Xiong F, Sun Z, Xu M. MVPainter: Accurate and detailed 3D texture generation via multi-view diffusion with geometric control. arXiv:2505.12635, 2025.

\bibitem{kynkanniemi2024limited}
Kynk\"a\"anniemi T, Aittala M, Karras T, et al. Applying guidance in a limited interval improves sample and distribution quality in diffusion models. \textit{Advances in Neural Information Processing Systems}, 2024.

\bibitem{perla2026survey}
Perla S R K, Zhang H, Mahdavi-Amiri A. Advances in neural 3D mesh texturing: A survey. \textit{Computer Graphics Forum}, 2026, e70392.

\bibitem{guzman2026adaptive}
Guzm\'an J E, Mould D, Paquette E. Adaptive multiresolution exemplar-based texture synthesis on animated fluids. \textit{Computers \& Graphics}, 2026, 134: 104490.

\bibitem{lv2025interpretable}
Lv X, Wu Z, Xu J, et al. Interpretable procedural material graph generation via diffusion models from reference images. \textit{The Visual Computer}, 2025, 41(13): 11195--11205.

\bibitem{zeng2024paint3d}
Zeng X, Chen X, Qi Z, et al. Paint3D: Paint Anything 3D with Lighting-Less Texture Diffusion Models. \textit{Proceedings of the IEEE/CVF Conference on Computer Vision and Pattern Recognition}, 2024.

\bibitem{cheng2024mvpaint}
Cheng W, Mu J, Zeng X, et al. MVPaint: Synchronized multi-view diffusion for painting anything 3D. \textit{Proceedings of the IEEE/CVF Conference on Computer Vision and Pattern Recognition}, 2025.

\bibitem{lai2025hunyuan3d}
Zhao Z, Lai Z, Lin Q, et al. Hunyuan3D 2.0: Scaling diffusion models for high resolution textured 3D assets generation. arXiv:2501.12202, 2025.

\bibitem{huang2024mvadapter}
Huang Z, et al. MV-Adapter: Multi-view consistent image generation made easy. \textit{Proceedings of the IEEE/CVF International Conference on Computer Vision}, 2025.

\bibitem{tang2023mvdiffusion}
Tang S, Zhang F, Chen J, Wang P, Furukawa Y. MVDiffusion: Enabling holistic multi-view image generation with correspondence-aware diffusion. \textit{Advances in Neural Information Processing Systems}, 2023.

\bibitem{liu2024synchronized}
Liu Y X, Xie M S, Liu H Y, Wong T T. Text-guided texturing by synchronized multi-view diffusion. \textit{SIGGRAPH Asia 2024 Conference Papers}, 2024.

\bibitem{xu2026wondertex}
Xu Q, Han X G, Zhang L. WonderTex: Consistent-and-seamless texture generation with text-guided multi-view image diffusion models. \textit{IEEE Transactions on Visualization and Computer Graphics}, 2026, 32(7): 5815--5825.

\bibitem{mou2023t2iadapter}
Mou C, Wang X, Xie L, et al. T2I-Adapter: Learning adapters to dig out more controllable ability for text-to-image diffusion models. \textit{Proceedings of the AAAI Conference on Artificial Intelligence}, 2024.

\bibitem{zhang2023controlnet}
Zhang L, Rao A, Agrawala M. Adding conditional control to text-to-image diffusion models. \textit{Proceedings of the IEEE/CVF International Conference on Computer Vision}, 2023.

\bibitem{ye2023ipadapter}
Ye H, Zhang J, Liu S, Han X, Yang W. IP-Adapter: Text compatible image prompt adapter for text-to-image diffusion models. arXiv:2308.06721, 2023.

\bibitem{hu2022lora}
Hu E J, Shen Y, Wallis P, et al. LoRA: Low-rank adaptation of large language models. \textit{International Conference on Learning Representations}, 2022.

\bibitem{liu2024dora}
Liu S Y, Wang C Y, Yin H, et al. DoRA: Weight-decomposed low-rank adaptation. \textit{International Conference on Machine Learning}, 2024.

\bibitem{zhang2023adalora}
Zhang Q, Zuo S, Liang C, et al. AdaLoRA: Adaptive budget allocation for parameter-efficient fine-tuning. \textit{International Conference on Learning Representations}, 2023.

\bibitem{kulkarni2026ssi}
Kulkarni A S. Scheduled Style Injection: Expanding the Style-Content Pareto Frontier in Training-Free Diffusion-Based Style Transfer. \textit{Proceedings of the IEEE/CVF Conference on Computer Vision and Pattern Recognition Workshops (CVPRW), NTIRE Workshop}, 2026.

\bibitem{ho2020ddpm}
Ho J, Jain A, Abbeel P. Denoising diffusion probabilistic models. \textit{Advances in Neural Information Processing Systems}, 2020.

\bibitem{song2021ddim}
Song J, Meng C, Ermon S. Denoising diffusion implicit models. \textit{International Conference on Learning Representations}, 2021.

\bibitem{ho2022cfg}
Ho J, Salimans T. Classifier-free diffusion guidance. \textit{NeurIPS Workshop on Deep Generative Models and Downstream Applications}, 2021. arXiv:2207.12598.

\bibitem{castillo2023adaptive}
Castillo A, Kohler J, P\'erez J C, et al. Adaptive guidance: Training-free acceleration of conditional diffusion models. \textit{Proceedings of the AAAI Conference on Artificial Intelligence}, 2025.

\bibitem{wang2004ssim}
Wang Z, Bovik A C, Sheikh H R, Simoncelli E P. Image quality assessment: From error visibility to structural similarity. \textit{IEEE Transactions on Image Processing}, 2004, 13(4): 600--612.

\bibitem{zhang2018lpips}
Zhang R, Isola P, Efros A A, et al. The unreasonable effectiveness of deep features as a perceptual metric. \textit{CVPR}, 2018.

\bibitem{sad2024apg}
Sadat S, Hilliges O, Weber R M. Eliminating oversaturation and artifacts of high guidance scales in diffusion models. \textit{International Conference on Learning Representations}, 2025.

\bibitem{chung2024styleid}
Chung J, Hyun S, Heo J P. Style Injection in Diffusion: A Training-free Approach for Adapting Large-scale Diffusion Models for Style Transfer. \textit{Proceedings of the IEEE/CVF Conference on Computer Vision and Pattern Recognition}, 2024.

\bibitem{controlnet2023layercontrol}
ControlNet project. Official depth-to-image demo: layer-dependent control scales. Source code, fixed commit f2e1bc38dc8d925fb542be4c138c46013ace1988, 2023. \url{https://github.com/lllyasviel/ControlNet/blob/f2e1bc38dc8d925fb542be4c138c46013ace1988/gradio_depth2image.py}. Accessed 12 October 2026.

\bibitem{diffusers2023controlnet}
Hugging Face. Diffusers v0.21.0: ControlNet conditioning scale and guidance start/end. Official versioned documentation, 2023. \url{https://huggingface.co/docs/diffusers/v0.21.0/api/pipelines/controlnet}. Accessed 12 October 2026.

\end{thebibliography}
\endgroup
\end{document}
